\section{Writing Guidelines}

\subsection{Abstract and Keywords}

The abstract should state the problem, the method, and the result in one compact paragraph. For a class or project paper, it is acceptable to mention that the work is an implementation, evaluation, replication, or analysis rather than a new method. Avoid writing an abstract that only introduces the topic. A reader should understand what was built, what was measured, and what conclusion the paper reaches.

Keywords should be specific enough to help classification. For example, a paper about model evaluation may use terms such as \textit{benchmarking}, \textit{classification}, \textit{latency}, and \textit{reproducibility}. A paper about a software system may use terms such as \textit{architecture}, \textit{data pipeline}, \textit{performance analysis}, and \textit{deployment}.

\subsection{Problem Statement}

A good problem statement should identify the gap or failure mode being addressed. Depending on the topic, the issue may be inaccurate predictions, high latency, poor usability, data loss, unreliable synchronization, weak scalability, deployment friction, or inconsistent measurement. The statement should then explain why the chosen method addresses that issue. This prevents the paper from becoming a general survey with no concrete claim.

\subsection{Design and Implementation}

Design sections should describe abstractions and goals. Implementation sections should describe concrete choices. For example, a design section may state that a system separates data ingestion from analysis and reporting. The implementation section should then name the modules, fields, command-line options, and data files used to realize that separation.

When presenting code-level work, avoid pasting long source listings. A short algorithm is useful when it clarifies control flow, but the repository should hold the complete implementation. The paper should explain decisions that are not obvious from code alone: why a parameter was selected, why a data format was chosen, what assumptions are made, and which cases are intentionally rejected.

\subsection{Evaluation}

Evaluation should connect back to the original claim. If the claim is performance, report runtime, memory use, bandwidth, or storage overhead. If the claim is correctness, report success and failure cases. If the claim is usability, report task completion, error rate, or qualitative feedback. Even a small experiment is stronger when it has explicit scenarios and expected outcomes.
